\documentclass[12pt]{article}
\usepackage[a4paper,margin=1in]{geometry}
\usepackage[utf8]{inputenc}
\usepackage[T1]{fontenc}
\usepackage{amsmath,amssymb}
\usepackage{textcomp}
\usepackage{parskip}
\begin{document}
\begin{center}
{\Large \textbf{1998 STEP III}}
\end{center}

\subsection*{Page 1}
\begin{flushleft}
\bigskip\textbf{Section A: Pure Mathematics}\par
1 Let\par
\[f(x) = sin x + 2 cos x + 1\]
for 0 6 x 6 2\pi. Sketch the curve y = f(x), giving the coordinates of the stationary points.\par
Now let\par
af(x) + b\par
\[g(x) = ad 6= bc , d 6= -3c , d 6= c .\]
cf(x) + d\par
Show that the stationary points of y = g(x) occur at the same values of x as those of y = f(x),\par
and find the corresponding values of g(x).\par
\[Explain why, if d/c < -3 or d/c > 1, |g(x)| cannot be arbitrarily large.\]
2 Let\par
\[I(a, b) = ta(1 - t)b dt (a > 0, b > 0).\]
\[\int\]
\textbf{(i)} Show that I(a, b) = I(b, a),\par
\textbf{(ii)} Show that I(a, b) = I(a + 1, b) + I(a, b + 1).\par
\textbf{(iii)} Show that (a+1)I(a, b) = bI(a+1, b-1) when a and b are positive and hence calculate\par
I(a, b) when a and b are positive integers.\par
\bigskip\textbf{Question 3.} The value V of a bond after N days is determined by the equation\par
N\par
\[V N+1 = (1 + c)V N - d (c > 0, d > 0),\]
\[where c and d are given constants. By looking for solutions of the form V = AkT + B for\]
T\par
some constants A, B and k, or otherwise, find V\par
N\par
\[in terms of V^0.\]
What is the solution for c = 0? Show that this is the limit (for fixed N ) as c \to 0 of your\par
\[solution for c > 0.\]
\bigskip\textbf{Question 4.} Show that the equation (in plane polar coordinates) r = cos θ, for - \pi 6 θ 6 \pi , represents a\par
2 2\par
circle.\par
Sketch the curve r = cos 2θ for 0 6 θ 6 2\pi, and describe the curves r = cos 2nθ, where n is\par
an integer. Show that the area enclosed by such a curve is independent of n.\par
\[Sketch also the curve r = cos 3θ for 0 6 θ 6 2\pi.\]
\end{flushleft}
\subsection*{Page 2}
\begin{flushleft}
\bigskip\textbf{Question 5.} The exponential of a square matrix A is defined to be\par
\infty\par
\[exp(A) = Ar ,\]
r!\par
\[r=0\]
X\par
\[where A^0 = I and I is the identity matrix.\]
Let\par
0 -1\par
\[M = .\]
1 0\par
µ ¶\par
Show that M2 = -I and hence express exp(θM) as a single 2 \times 2 matrix, where θ is a real\par
number. Explain the geometrical significance of exp(θM).\par
Let\par
0 1\par
\[N = .\]
0 0\par
µ ¶\par
Express similarly exp(sN), where s is a real number, and explain the geometrical significance\par
of exp(sN).\par
For which values of θ does\par
\[exp(sN) exp(θM) = exp(θM) exp(sN)\]
for all s? Interpret this fact geometrically.\par
\bigskip\textbf{Question 6.} (i) Show that four vertices of a cube, no two of which are adjacent, form the vertices of\par
a regular tetrahedron. Hence, or otherwise, find the volume of a regular tetrahedron\par
whose edges are of unit length.\par
\textbf{(ii)} Find the volume of a regular octahedron whose edges are of unit length.\par
\textbf{(iii)} Show that the centres of the faces of a cube form the vertices of a regular octahedron.\par
Show that its volume is half that of the tetrahedron whose vertices are the vertices of\par
the cube.\par
[A regular tetrahedron (octahedron) has four (eight) faces, all equilateral triangles.]\par
\end{flushleft}
\subsection*{Page 3}
\begin{flushleft}
\bigskip\textbf{Question 7.} Sketch the graph of f(s) = es(s - 3) + 3 for 0 6 s < \infty. Taking e \approx 2.7, find the smallest\par
\[positive integer, m, such that f(m) > 0.\]
Now let\par
\[x^3\]
\[b(x) =\]
ex/T - 1\par
where T is a positive constant. Show that b(x) has a single turning point in 0 < x < \infty. By\par
\[considering the behaviour for small x and for large x, sketch b(x) for 0 6 x < \infty.\]
Let\par
\infty\par
\[b(x) dx = B,\]
\[\int\]
\[which may be assumed to be finite. Show that B = KT n where K is a constant, and n is an\]
integer which you should determine.\par
T m\par
Given that B \approx 2 b(x)dx, use your graph of b(x) to find a rough estimate for K.\par
R\par
\bigskip\textbf{Question 8.} (i) Show that the line r = b + λm, where m is a unit vector, intersects the sphere r.r = a2\par
at two points if\par
\[a 2 > b.b - (b.m) 2 .\]
Write down the corresponding condition for there to be precisely one point of intersec-\par
\[tion. If this point has position vector p, show that m.p = 0.\]
\textbf{(ii)} Now consider a second sphere of radius a and a plane perpendicular to a unit vector n.\par
The centre of the sphere has position vector d and the minimum distance from the\par
origin to the plane is l. What is the condition for the plane to be tangential to this\par
second sphere?\par
\textbf{(iii)} Show that the first and second spheres intersect at right angles (i.e. the two radii to\par
each point of intersection are perpendicular) if\par
\[d.d = 2a 2 .\]
\end{flushleft}
\subsection*{Page 4}
\begin{flushleft}
\bigskip\textbf{Section B: Mechanics}\par
\bigskip\textbf{Question 9.} A uniform right circular cone of mass m has base of radius a and perpendicular height h\par
from base to apex. Show that its moment of inertia about its axis is 3 ma2, and calculate its\par
moment of inertia about an axis through its apex parallel to its base.\par
[Any theorems used should be stated clearly.]\par
The cone is now suspended from its apex and allowed to perform small oscillations. Show that\par
their period is\par
\[4h2 + a^2\]
\[2\pi .\]
s 5gh\par
[You may assume that the centre of mass of the cone is a distance h from its apex.]\par
\bigskip\textbf{Question 10.} Two identical spherical balls, moving on a horizontal, smooth table, collide in such a way that\par
both momentum and kinetic energy are conserved. Let v 1 and v 2 be the velocities of the balls\par
0 0\par
before the collision and let v and v be the velocities of the balls after the collision, where\par
1 2\par
0 0\par
v 1, v 2, v\par
and v\par
are two-dimensional vectors. Write down the equations for conservation of\par
momentum and kinetic energy in terms of these vectors. Hence show that their relative speed\par
is also conserved.\par
Show that, if one ball is initially at rest but after the collision both balls are moving, their\par
final velocities are perpendicular.\par
Now suppose that one ball is initially at rest, and the second is moving with speed V . After\par
a collision in which they lose a proportion k of their original kinetic energy (0 6 k 6 1), the\par
direction of motion of the second ball has changed by an angle θ. Find a quadratic equation\par
satisfied by the final speed of the second ball, with coefficients depending on k, V and θ.\par
Hence show that k 6 1 .\par
\end{flushleft}
\subsection*{Page 5}
\begin{flushleft}
\bigskip\textbf{Question 11.} Consider a simple pendulum of length l and angular displacement θ, which is not assumed to\par
be small. Show that\par
1 dθ\par
\[l = g(cos θ - cos γ) ,\]
2 dt\par
µ ¶\par
where γ is the maximum value of θ. Show also that the period P is given by\par
\[P = 2 l γ sin 2 (γ/2) - sin 2 (θ/2) - 1 2 dθ .\]
s g 0\par
\[\int\]
¡ ¢\par
By using the substitution sin(θ/2) = sin(γ/2) sin φ, and then finding an approximate expres-\par
sion for the integrand using the binomial expansion, show that for small values of γ the period\par
is approximately\par
l γ2\par
\[2\pi 1 + .\]
s g 16\par
µ ¶\par
\end{flushleft}
\subsection*{Page 6}
\begin{flushleft}
\bigskip\textbf{Section C: Probability and statistics}\par
\bigskip\textbf{Question 12.} The mountain villages A, B, C and D lie at the vertices of a tetrahedron, and each pair of\par
villages is joined by a road. After a snowfall the probability that any road is blocked is p, and\par
is independent of the conditions of any other road. The probability that, after a snowfall, it\par
is possible to travel from any village to any other village by some route is P . Show that\par
2 3 2\par
\[P = 1 - p (6p - 12p + 3p + 4).\]
13 Write down the probability of obtaining k heads in n tosses of a fair coin. Now suppose that\par
k is known but n is unknown. A maximum likelihood estimator (MLE) of n is defined to be a\par
value (which must be an integer) of n which maximizes the probability of k heads.\par
A friend has thrown a fair coin a number of times. She tells you that she has observed one\par
head. Show that in this case there are two MLEs of the number of tosses she has made.\par
She now tells you that in a repeat of the exercise she has observed k heads. Find the two\par
MLEs of the number of tosses she has made.\par
She next uses a coin biased with probability p (known) of showing a head, and again tells you\par
that she has observed k heads. Find the MLEs of the number of tosses made. What is the\par
condition for the MLE to be unique?\par
\bigskip\textbf{Question 14.} A hostile naval power possesses a large, unknown number N of submarines. Interception of\par
radio signals yields a small number n of their identification numbers X (i = 1, 2, ..., n), which\par
i\par
are taken to be independent and uniformly distributed over the continuous range from 0 to\par
\[N . Show that Z^1 and Z^2, defined by\]
n\par
n + 1 2\par
\[Z^1 = Max\{X^1, X^2, ..., X\]
n\par
\[\} and Z^2 = X\]
i\par
,\par
n n\par
\[i=1\]
X\par
both have means equal to N .\par
\[Calculate the variance of Z^1 and of Z^2. Which estimator do you prefer, and why?\]
\end{flushleft}
\end{document}
